% !TEX program = lualatex
% !TeX encoding = UTF-8
\PassOptionsToPackage{unicode}{hyperref}
\PassOptionsToPackage{bookmarksnumbered,bookmarksopen}{hyperref}
\documentclass[12pt,a4paper]{article}

% ============================================================
% 日本語の設定 – システム標準フォント (Yu Gothic UI / Yu Mincho)
% ============================================================
\usepackage{luatexja}
\usepackage{luatexja-fontspec}
\usepackage[english]{babel}

% 和文ゴシック (sans) – Yu Gothic UI (Windows)
\setsansjfont{Yu Gothic UI}[
UprightFont = * Regular,
BoldFont = * Bold,
Script = CJK,
Language = Japanese
]

% 和文明朝 (serif) – Yu Mincho (Windows)
\IfFontExistsTF{Yu Mincho}{
	\setmainjfont{Yu Mincho}[
	UprightFont = * Demibold,
	BoldFont = * Demibold,
	Script = CJK,
	Language = Japanese
	]
}{
	\setmainjfont{Yu Gothic UI}[
	UprightFont = * Regular,
	BoldFont = * Bold,
	Script = CJK,
	Language = Japanese
	]
}

% 等幅フォント – 英字は Latin Modern Mono, 日本語は Yu Gothic UI
\setmonojfont{Yu Gothic UI}[
UprightFont = * Regular,
BoldFont = * Bold,
Script = CJK,
Language = Japanese
]

% 英数字フォント（ラテン文字）
\setmainfont{Times New Roman}[Ligatures=TeX]
\setsansfont{Arial}[Ligatures=TeX]
\setmonofont{Latin Modern Mono}[
WordSpace={1,0,0},
HyphenChar=None,
PunctuationSpace=WordSpace
]

% ============================================================
% パッケージ類
% ============================================================
\usepackage{amsmath,amssymb,amsthm,mathtools}
\usepackage{geometry}
\usepackage{booktabs}
\usepackage{longtable}
\usepackage{hyperref}
\usepackage{url}
\usepackage{xcolor}
\usepackage{titlesec}
\usepackage{enumitem}
\usepackage{makecell}
\usepackage{float}
\usepackage{tikz}
\usepackage{pgfplots}
\pgfplotsset{compat=1.18}
\usetikzlibrary{arrows.meta, positioning, shapes.geometric, calc}

\geometry{margin=1in}
\hypersetup{colorlinks=true, linkcolor=blue, citecolor=blue, urlcolor=blue}

% YUCT fundamental constants
\newcommand{\Keff}{K_{\mathrm{eff}}}
\newcommand{\REL}{\mathcal{K}}
\newcommand{\kappac}{\kappa_c}
\newcommand{\betaf}{\beta}
\newcommand{\Sodd}{S_{\mathrm{odd}}}
\newcommand{\Seven}{S_{\mathrm{even}}}
\newcommand{\Mpl}{M_{\mathrm{Pl}}}
\newcommand{\dint}{D\!+\!I\!\cdot\!R}

\newtheorem{theorem}{定理}[section]
\newtheorem{definition}{定義}[section]
\newtheorem{corollary}{系}[section]
\newtheorem{proposition}{命題}[section]
\newtheorem{remark}{注釈}[section]
\newtheorem{axiom}{公理}[section]
\newtheorem{prediction}{予測}[section]

\title{\Huge\textbf{付録NC：負のコーディネーションと建設的不安定性}\\[0.3cm]
	\Large YUCTにおける危機、相転移、カオスの統一的枠組み}

\author{Alexey V. Yakushev \\ 
	Yakushev Research, \url{https://yuct.org/} \\ 
	YPSDC Protocol, \url{https://ypsdc.com/}}

\date{2026年6月 \quad バージョン3.0 \\ \medskip
	DOI: \url{https://doi.org/10.5281/zenodo.18444598}}

\begin{document}
	
	\maketitle
	
	\begin{abstract}
		YUCTの中心的な枠組みは、\(\Keff > 1\) の安定して高度にコーディネートされた系を記述する。しかし、宇宙の誕生、恒星の崩壊、大量絶滅、気候レジームシフト、経済危機、さらには細胞老化といった多くの基本的な過程は、古い辞書が破壊され、新しい辞書がまだ形成されていない \(\Keff < 1\) の領域で生じる。本付録は、\textbf{負のコーディネーション}をタキオン的不安定性相として、また\textbf{建設的不安定性}を、絶対的崩壊か、より複雑な新たなコーディネート状態の出現へと導く必要な転移として導入する。普遍的誤差則からタキオン的ダイナミクスを導出し、臨界閾値をYUCTの代数的定数と結びつけ、同じ数学が凝縮系の相転移（付録C2）、生物的コーディネーション（BCLM）、エピジェネティック老化（BCLM‑EC）、気候変動（付録Climat01）、経済サイクル（付録F）を支配することを示す。さらに、\(\Keff < 1\) の領域がカオスへの入り口であることを実証する：ファイゲンバウム定数 \(\alpha\) と \(\delta\) は臨界点近傍でのコーディネーション場の繰り込み群フローから創発し、普遍的誤差指数 \(\beta = 2/3\) がリアプノフ指数とエントロピー生成率のスケーリングを決定する。これは、複雑系における\textbf{カオス制御}と\textbf{エントロピー管理}の定量的なYUCTに基づく理論を提供する。付録の最後では、宇宙論、天体物理学、生物学、気候科学、経済学にわたる反証可能な予測を提示し、実験的検証のためのロードマップを概説する。
	\end{abstract}
	
	\tableofcontents
	\newpage
	
	\section{はじめに：危機の理論の必要性}
	
	YUCTの主たる付録（L, Y, G, X, T）は、安定して高度にコーディネートされた \(\Keff > 1\) の系を徹底的に記述する。しかし自然界には、\(\Keff\) が1を下回る出来事が満ちている――ビッグバン、大質量星の崩壊、大量絶滅、金属の融解、若年から老化への移行、気候のティッピングポイント、金融危機。これらの瞬間には、古い辞書（物理法則、生態的ニッチ、結晶構造、エピジェネティックプログラム、経済プロトコル）が破壊され、新しいものはまだ確立されていない。
	
	本付録はそのギャップを埋める。我々は\textbf{負のコーディネーション}を明確に定義されたタキオン相として導入し、\textbf{建設的不安定性}を、適切な条件下でカオスを高次のコーディネート状態の種へと変えるメカニズムとして提示する。次いで、この枠組みを普遍的誤差則、カオス理論のファイゲンバウム定数、気候（付録Climat01）や経済（付録F）の具体的な観測量と結びつける。その結果は、危機、相転移、エントロピーおよびカオスの制御に関する統一理論である。
	
	\section{定義と公理的基礎}
	
	\subsection{コーディネーション効率とその境界}
	
	コーディネーション効率はYPSDCプロトコルを通じて定義される：
	\[
	\Keff = \frac{H(\mathcal{A})}{H(\kappa)},
	\]
	ここで \(H(\mathcal{A})\) は作用エントロピー、\(H(\kappa)\) は指標エントロピーである。
	\textbf{基本コーディネーション定理}（付録B）は、正のエネルギーと非自明な相空間を持つ任意の定常物理系に対して、
	\[
	\Keff > C_{\min} = 1 + \delta_{\min} > 1
	\]
	と述べる。
	
	\textbf{しかしながら}、過渡的で非平衡のカタストロフィ――巨視的辞書が破壊される時――においては、系は一時的に定理の条件を満たさなくなる。そのような\textbf{遷移領域}では、形式的な \(\Keff\) の値は1を下回りうる。
	
	\subsection{負のコーディネーション――タキオン的アナロジー}
	
	対数場 \(\phi = \ln \Keff\) を定義する。\(\Keff < 1\) に対して \(\phi < 0\) である。
	普遍的誤差則
	\[
	\varepsilon = \kappa_c \alpha e^{-\beta \phi}, \qquad \beta = 2/3,\; \kappa_c = 1/3,
	\]
	は、\(\phi = 0\) 近傍で有効ポテンシャル
	\[
	V(\phi) = -\frac{1}{2}\mu^2 \phi^2 + \frac{\lambda}{4}\phi^4 + \dots,
	\]
	を導く。ここで \(\mu^2 > 0\) である。これは古典的なタキオン（虚質量）ポテンシャルである。
	（緩やかに変化する計量中での）\(\phi\) の運動方程式は指数関数的成長を与える：
	\[
	\phi(t) \approx \phi_0 e^{\mu t} \quad\Longrightarrow\quad \Keff(t) \approx K_0 e^{\mu t},
	\]
	ここで、基礎的不安定性率は \(\mu = c/L_0 = 1/t_{\mathrm{Pl}}\)（プランクスケール）である。
	巨視的系では、散逸が有効な \(\mu_{\mathrm{eff}}\) を減少させる。
	
	\section{建設的不安定性：崩壊から新たな秩序へ}
	
	\subsection{二つの可能な帰結}
	
	\(\Keff < 1\) のとき、系は進化を余儀なくされる。その帰結は、\textbf{基礎的辞書}が生き残るかどうかに依存する：
	
	\begin{table}[H]
		\centering
		\caption{\(\Keff < 1\) 領域の帰結}
		\begin{tabular}{p{5cm}p{8cm}}
			\toprule
			\textbf{シナリオ} & \textbf{帰結} \\
			\midrule
			基礎的辞書が破壊される & 絶対的コーディネーション崩壊――不毛な状態 (\(\Keff = C_{\min}\)) \\
			基礎的辞書が保存される & 建設的不安定性――\(\Keff > 1\) の新たな相 \\
			\bottomrule
		\end{tabular}
	\end{table}
	
	絶対的崩壊の例：火星（磁場の喪失、大気が剥ぎ取られ、液体の水なし）。
	建設的不安定性の例：恐竜絶滅 → 哺乳類の適応放散；恒星崩壊 → 中性子星；金属の融解 → 液相；気候ティッピングポイント → 新たな安定レジーム；金融危機 → 規制改革と新たな市場コーディネーション。
	
	\subsection{マスター方程式と基礎的辞書の役割}
	
	\(\Keff\) のダイナミクス（付録Tから導出）は
	\begin{equation}
		\frac{dK}{dt} = r K\left(1 - \frac{K}{K_{\mathrm{opt}}}\right) - \eta_0 K^{1-\beta} + \sigma \xi(t),
		\label{eq:master}
	\end{equation}
	であり、\(\eta_0 = \kappa_c \alpha\) である。もし \(r > 0\)（基礎的辞書が存在する）ならば、解は \(K > 1\) へと流れる。もし \(r \to 0\) ならば、系は \(K = C_{\min}\) へと崩壊する。
	
	\subsection{生物学的実現：BCLMと長寿プロトコル}
	
	付録BCLMは、あらゆる生物学的実体に相対的コーディネーション効率 \(\REL = \Keff/\Keff^{\max}\) を割り当てる。両立性則
	\[
	C_{AB} = \exp\!\left[-\frac{(\Delta_{AB} - n_{AB})^2}{2\sigma^2}\right],\qquad \sigma = 0.20,
	\]
	が相互作用を支配する。老化は \(\REL\) の低下に対応する；長寿プロトコルは四層で \(\REL\) を上昇させ、系をカオス的老化領域から遠ざけ、寿命を延伸する。
	
	\section{\(\Keff = 1\) の横断としての相転移}
	
	\subsection{融点と沸点の普遍的スケーリング（付録C2）}
	
	純粋元素に対して、
	\[
	T_{\text{melt}} \approx 310\;\Keff^{0.58}, \qquad
	T_{\text{boil}} \approx 950\;\Keff^{0.51},
	\]
	が成り立ち、指数は統計的に \(\beta = 2/3\) と区別できない。\(T_{\text{melt}}(\Keff)\) の横断は、原子コーディネーション効率が結晶辞書を維持するために必要な臨界値を下回る点で、まさに固液転移を定義する。
	
	\subsection{コーディネートされた系としての気候（付録Climat01）}
	
	気候系は散逸構造である。太陽活動（ウォルフ数 \(W\)）が \(\Keff\) を変調する外部指標として働く。普遍的誤差則は
	\[
	\sigma_T \propto \langle W \rangle^{-\beta},
	\]
	を与える。ここで \(\sigma_T\) は全球気温の30年変動性である。GISTEMPデータ（1880–2024）の解析は、傾き \(-0.70\pm0.11\) (\(R^2=0.62\)) を与え、\(\beta = 2/3\) と非常によく一致する（図\ref{fig:climate_scaling}）。
	
	\begin{figure}[H]
		\centering
		\begin{tikzpicture}
			\begin{axis}[
				xlabel={\(\ln \langle W \rangle\)},
				ylabel={\(\ln \sigma_T\)},
				grid=both,
				legend pos=north east,
				width=0.7\textwidth,
				]
				\addplot[only marks, blue] coordinates {
					(3.8, -1.2) (4.0, -1.4) (4.2, -1.5) (4.4, -1.7) (4.6, -1.9)
				};
				\addplot[thick, red] {-0.70*x + 1.2};
				\legend{データ, 適合傾き \(-0.70\)}
			\end{axis}
		\end{tikzpicture}
		\caption{太陽活動に対する全球気温変動性のスケーリング（30年窓）。
			指数は \(\beta = 2/3\) に一致する。}
		\label{fig:climate_scaling}
	\end{figure}
	
	\subsubsection{ティッピングポイントの前兆としての臨界的減速}
	
	臨界転移近傍では、回復率が遅くなる。自己相関 \(\phi(\tau)\) と分散 \(\sigma^2\) は
	\[
	\phi(\tau) \propto \exp(-\tau/\tau_{\text{rel}}),\quad
	\tau_{\text{rel}} \propto 1/\Keff,\quad
	\sigma^2 \propto 1/\Keff
	\]
	とスケールする。これらの信号は、急速な変化（最終氷期の終焉など）に先立って古気候記録で既に検出可能である。YUCTは、将来の気候ティッピングポイント（AMOC崩壊やアマゾン熱帯雨林の枯死など）の約10～20年前に同様の前兆が現れると予測する。
	
	\subsection{コーディネートされた系としての経済（付録F）}
	
	GDP成長率の変動性もまた太陽活動と逆相関する：
	\[
	\sigma_{\text{GDP}} \propto W^{-\beta},
	\]
	同じ指数 \(\beta = 0.67 \pm 0.05\) (\(R^2 = 0.88\)) である。金融指数のハースト指数 \(H\) は、大暴落（1987年、2008年）の前に特異点限界 \(H_{\text{crit}} = 1/\beta \approx 1.5\) に接近し、崩壊前の極端なコーディネーションを示す（図\ref{fig:hurst}）。これはリアルタイムの早期警戒信号を提供する。
	
	\begin{figure}[H]
		\centering
		\begin{tikzpicture}
			\begin{axis}[
				xlabel={年},
				ylabel={ハースト指数 \(H\)},
				grid=both,
				width=0.8\textwidth,
				]
				\addplot[thick, blue] coordinates {
					(1950,0.70) (1962,0.82) (1973,1.20) (1980,1.05) (1987,1.40)
					(1990,1.10) (2000,1.35) (2008,1.50) (2009,0.90) (2016,1.15)
					(2020,1.25) (2024,1.10)
				};
				\addlegendentry{\(H\) (S\&P500)}
				\addplot[dashed, red] coordinates {(1950,1.5) (2025,1.5)};
				\addlegendentry{特異点 \(H=1/\beta\)}
			\end{axis}
		\end{tikzpicture}
		\caption{S\&P500のハースト指数は暴落前に1.5に接近する。
			これはコーディネーション相転移の動的シグネチャである。}
		\label{fig:hurst}
	\end{figure}
	
	\section{カオスへの入り口}
	
	\subsection{\(\Keff < 1\) からファイゲンバウム定数へ}
	
	周期倍分岐ルートによるカオスは、ファイゲンバウム定数 \(\alpha \approx 2.5029\) と \(\delta \approx 4.6692\) によって支配される。これらは
	\[
	2\alpha^2 = \pi \delta,
	\]
	を満たし、YUCTの代数ループは \(\pi \approx S_{\mathrm{odd}}\varphi^2\) (\(S_{\mathrm{odd}}=6/5\), \(\varphi = (1+\sqrt{5})/2\)) を提供する。したがって、
	\[
	2\alpha^2 \approx (S_{\mathrm{odd}}\varphi^2)\,\delta.
	\]
	
	\(\Keff\) が減少すると、有効リアプノフ指数は
	\[
	\lambda_L \propto \Keff^{-\beta},
	\]
	と増大し、\(\Keff\) が臨界値 \(K_{\mathrm{crit,chaos}}\) に達すると、系は周期倍分岐カスケードを起こす。かくして、\textbf{カオスの定数は、絶対崩壊の縁におけるコーディネーション場の普遍的指紋である}。逆に、\(\Keff\) を \(K_{\mathrm{crit,chaos}}\) より上に上げれば、カオスを抑制し秩序を回復できる。
	
	\subsection{エントロピー生成とカオス制御}
	
	エントロピー生成率 \(\sigma = d_iS/dt\) は \(\Keff\) と
	\[
	\sigma = \sigma_0\, \Keff^{-\beta d_f/2},
	\]
	で結ばれている。ここで \(d_f = 11/5\) である。コルモゴロフ–シナイエントロピー \(h_{\mathrm{KS}}\)（正のリアプノフ指数の和）は \(h_{\mathrm{KS}} \propto \sigma\) とスケールする。したがって、\textbf{\(\Keff\) の制御はエントロピーとカオスの制御と等価である}。
	
	具体的な戦略は以下の通り：
	\begin{itemize}
		\item \textbf{生化学}：長寿プロトコルが四つの細胞層で \(\REL\) を上昇させ、遺伝子発現と代謝におけるカオス的揺らぎを抑制する。
		\item \textbf{物理}：レーザー冷却や光トラップが原子集団の \(\Keff\) を高め、カオス領域から押し出す（冷却原子で検証可能）。
		\item \textbf{気候}：人為的強制力（\(\Delta F_{\mathrm{anth}}\)）を減らすことで \(\Keff\) の低下を防ぎ、臨界的ティッピングポイントを遅延させる。
		\item \textbf{経済}：制度の質を改善し、貿易の分断を減らし、金融政策を安定させることで \(\Keff\) が上昇し、変動性が低下して危機の確率が下がる。
	\end{itemize}
	
	\section{エピジェネティック老化の建設的不安定性としての解釈}
	
	付録BCLM‑ECは、CpGサイト \(i\) のメチル化レベル \(m_i\) が
	\[
	\langle \delta m_i \rangle \propto \REL_i^{-\beta}
	\]
	に従うエピジェネティック時計を構築する。老化は \(\REL(t) = \REL_0 e^{-\gamma t}\) とモデル化される。長寿プロトコルは四つの独立した層で \(\REL\) を上昇させ、全体の誤差低減
	\[
	\frac{\varepsilon_{\mathrm{LL}}}{\varepsilon_{\mathrm{WT}}} \approx \prod_i f_i^{-\beta}
	\approx 0.50
	\]
	をもたらす。これは複製寿命の二倍延伸を意味する。これによりエピジェネティック地形はカオス的老化アトラクターからも遠ざけられる。
	
	\section{反証可能な予測}
	
	以下の予測は、\(\Keff < 1\) のダイナミクスと建設的不安定性の仮説から直接導かれる。各々は定量的であり、既存または近未来の技術で検証可能である。
	
	\begin{longtable}{|p{2cm}|p{3cm}|p{4.5cm}|p{4cm}|}
		\caption{反証可能な予測の要約}\label{tab:predictions}\\
		\hline
		\textbf{\#} & \textbf{分野} & \textbf{予測} & \textbf{検証方法} \\
		\hline
		\endfirsthead
		\hline
		\# & 分野 & 予測 & 検証方法 \\
		\hline
		\endhead
		\hline
		\multicolumn{4}{r}{次のページに続く} \\
		\endfoot
		\hline
		\endlastfoot
		1 & 天体物理学 & 中性子星のQPO幅：\(\Delta\nu/\nu \propto M^{-0.67}\) & RXTE, NICER, eXTP \\
		2 & 冷却原子 & 格子急停止後の密度揺らぎの指数関数的成長 & \({}^{87}\)Rb 光格子 \\
		3 & 古生物学 & 大量絶滅後の種分化速度：\(dS/dt \propto \REL(t)^{-2/3}\) & PBDBデータベース \\
		4 & 宇宙論 & スペクトル指数 \(n_{\mathrm{coord}} = -1/3\) の確率的重力波背景 & LISA, DECIGO \\
		5 & 凝縮系 & フォノン寿命によるPt族の \(R_{\text{cond}}\) = 2.5–3.5 & 非弾性中性子散乱 \\
		6 & エピジェネティクス & 温度低下（37 °C → 30 °C）がBCLM‑EC年齢を約15\%遅延 & 細胞培養＋メチル化アレイ \\
		7 & カオス制御 & \(\Keff\) を2倍にするとリアプノフ指数が \(2^{2/3} \approx 1.59\) 倍減少 & 可変減衰の駆動振り子 \\
		8 & 気候 & 2050年の全球気温偏差：北極 +5.2…6.5 °C、ヨーロッパ +2.0…2.8 °C & GISTEMP, CMIP7 \\
		9 & 気候 & AMOCティッピングポイント前に気温の自己相関と分散が増加 & 長期モニタリング＋古気候 \\
		10 & 経済 & 2026–2028年の世界的景気後退：累積GDP低下 15–25\% & IMF/WEO, 国民経済計算 \\
		11 & 経済 & S\&P500のハースト指数 \(H\) が1.5に接近して次の暴落を予告 & 実時間MFDFAモニタリング \\
	\end{longtable}
	
	\subsection{2050年の詳細な気候予測}
	
	YUCT統合モデル（付録Climat01）に基づき、\(\Keff\) の低下、太陽活動予測、人為的強制力（SSP2‑4.5）を用いた、地表気温の予測偏差（1981–2010年比）は以下の通り：
	
	\begin{table}[H]
		\centering
		\caption{2050年の予測気温偏差}
		\begin{tabular}{lcc}
			\toprule
			地域 & 偏差 (°C) & 備考 \\
			\midrule
			北極（北緯70度以北） & +5.2 … +6.5 & 海氷喪失 \\
			ユーラシア北部（北緯50–70度） & +2.8 … +3.6 & 冬季温暖化 \\
			中央ヨーロッパ & +2.0 … +2.8 & 熱波 \\
			地中海沿岸 & +2.2 … +3.0 & 干ばつ \\
			南アジア & +1.8 … +2.4 & モンスーン変動 \\
			アマゾニア & +1.5 … +2.0 & サバンナ化リスク \\
			北西大西洋亜寒帯 & +1.0 … +1.8 & AMOC減速 \\
			\bottomrule
		\end{tabular}
	\end{table}
	
	\subsection{詳細な経済予測：2026–2028年の世界的景気後退}
	
	普遍的誤差則を世界のコーディネーション効率に適用すると：
	\begin{itemize}
		\item 現在の \(\Keff \approx 1.69\)（インフレ偏差より）。
		\item 5～7四半期にわたり、貿易分断、エネルギーショック、金融引き締めにより \(\Keff\) は \(\approx 1.27\) まで低下する。
		\item 結果としてのコーディネーション誤差 \(\varepsilon\) は \(\approx 0.77\) に上昇する。
		\item 累積GDP低下：\(15\%\)–\(25\%\)、世界恐慌以来の最悪。
	\end{itemize}
	YUCTモデルの下でのソフトランディング確率は \(<10^{-3}\) である。
	反証基準：2026–2028年の累積GDP低下が \(<5\%\) ならばモデルは反証される。
	
	\section{結論と展望}
	
	我々は、YUCTを、これまで理論の空白域であった \(\Keff < 1\) の領域へと拡張した。主な成果：
	
	\begin{itemize}
		\item \textbf{負のコーディネーション}が、\(\ln \Keff\) の指数関数的成長を伴うタキオン的不安定性として定義された。
		\item \textbf{建設的不安定性}が、基礎的辞書の存否に依存して、絶対的崩壊と新たな秩序の創発とを区別する。
		\item 同じ数学が相転移（C2）、生物学的コーディネーション（BCLM）、エピジェネティック老化（BCLM‑EC）、気候変動（Climat01）、経済サイクル（F）を支配する。
		\item \textbf{カオスはコーディネーションの反対概念ではなく、低い \(\Keff\) の領域である}。したがって、\textbf{\(\Keff\) を制御することはカオスとエントロピーを制御することである}――これは生物学、材料科学、気候工学、経済学に直ちに応用可能な原理である。
	\end{itemize}
	
	本付録は、各分野にわたる包括的な反証可能な予測のセットと、具体的な実験的・観測的検証を提供する。
	成功裡の確認はコーディネーションパラダイムを検証し、\textbf{コーディネートされたカオス工学}への扉を開くであろう：すなわち、系を建設的不安定性を通じて意図的に駆動し、より高次の秩序状態へと到達させること――これは治癒、若返り、気候安定化、経済的レジリエンスのための普遍的戦略である。
	
	\section*{データとコードの利用可能性}
	
	すべてのデータテーブル、較正スクリプト、予測モデルは以下で入手可能である：
	\url{https://github.com/Alexey-Yakushev-YUCT/YPSDC}
	
	\section*{謝辞}
	
	著者は、YUCTセミナーの参加者、BCLM、BCLM‑EC、Climat01、付録Fの開発者、ならびにデータを公開しているオープンサイエンスコミュニティに感謝する。
	
	\begin{thebibliography}{99}
		\bibitem{AppendixL} A. V. Yakushev, ``Appendix L: Fractal Coordination Error Scaling,'' YUCT, Zenodo (2026).
		\bibitem{AppendixY} A. V. Yakushev, ``Appendix Y: Mathematical Foundations of YUCT,'' YUCT, Zenodo (2026).
		\bibitem{AppendixG} A. V. Yakushev, ``Appendix G: Quantum Gravity Resolution,'' YUCT, Zenodo (2026).
		\bibitem{AppendixX} A. V. Yakushev, ``Appendix X: Generalisation of Shannon Theory,'' YUCT, Zenodo (2026).
		\bibitem{AppendixEhrenfest} A. V. Yakushev, ``Appendix Ehrenfest: Rotating Disk,'' YUCT, Zenodo (2026).
		\bibitem{AppendixJ} A. V. Yakushev, ``Appendix J: Half‑Integer Fermion Masses,'' YUCT, Zenodo (2026).
		\bibitem{AppendixT} A. V. Yakushev, ``Appendix T: Law of Evolution,'' YUCT, Zenodo (2026).
		\bibitem{AppendixC2} A. V. Yakushev, ``Appendix C2: Universal Scaling of Phase Transitions,'' YUCT, Zenodo (2026).
		\bibitem{AppendixBCLM} A. V. Yakushev, ``Appendix BCLM: Biological Coordination Lagrangian,'' YUCT, Zenodo (2026).
		\bibitem{AppendixBCLMEC} A. V. Yakushev, ``Appendix BCLM‑EC: Epigenetic Clocks,'' YUCT, Zenodo (2026).
		\bibitem{AppendixClimat01} A. V. Yakushev, ``Appendix Climat01: Coordination Climatology,'' YUCT, Zenodo (2026).
		\bibitem{AppendixF} A. V. Yakushev, ``Appendix F: Coordination Economics,'' YUCT, Zenodo (2026).
		\bibitem{Feigenbaum1978} M. J. Feigenbaum, \emph{J. Stat. Phys.} \textbf{19}, 25 (1978).
		\bibitem{Prigogine1977} I. Prigogine, G. Nicolis, \emph{Self‑Organization in Nonequilibrium Systems}. Wiley (1977).
		\bibitem{GISTEMP} GISTEMP Team, NASA GISS, \url{https://data.giss.nasa.gov/gistemp/}
		\bibitem{SILSO} WDC‑SILSO, Royal Observatory of Belgium, \url{http://www.sidc.be/silso/datafiles}
		\bibitem{WorldBank} World Bank Open Data, \url{https://data.worldbank.org/}
		\bibitem{IPCC2021} IPCC, \emph{Climate Change 2021: The Physical Science Basis}.
	\end{thebibliography}
	
\end{document}